\documentclass[11pt]{article}
\usepackage[margin=1in]{geometry}
\usepackage{amsmath,amssymb,amsthm}
\usepackage{graphicx}
\usepackage{hyperref}

\newtheorem{theorem}{Theorem}
\newtheorem{proposition}{Proposition}
\newtheorem{lemma}{Lemma}
\newtheorem{remark}{Remark}

\newcommand{\E}{\mathbb E}
\newcommand{\Pp}{\mathbb P}
\newcommand{\R}{\mathbb R}
\newcommand{\cB}{\mathcal B}
\newcommand{\norm}[1]{\left\lVert #1\right\rVert}
\newcommand{\Lq}{L^q(S)}
\newcommand{\sfq}{\mathsf{SF}_q}
\newcommand{\supp}{\operatorname{supp}}

\title{The \(L^q\), \(p=1\) Endpoint for \(1<q\le 2\)}
\author{Agent \texttt{agent\_lane\_19}}
\date{2026-06-16}

\begin{document}
\maketitle

\section*{Status}

This packet gives a candidate partial solution to Remark 12.10 of
Yaroslavtsev's paper \emph{Local characteristics and tangency of
vector-valued martingales} \cite{Yaroslavtsev2019}.  The source asks whether
characteristic domination implies maximal \(L^p\) comparison for general
local martingales.  Earlier packets in this run prove the \(L^q(S)\)-valued
case for \(1<p,q<\infty\), the Hilbert case for all finite \(p\), and a
conditional reduction of the remaining \(L^q(S)\), \(p=1\) endpoint to one
square-function comparison.

The new result here proves that square-function comparison for \(1<q\le 2\).
Consequently the full \(L^q(S)\)-valued characteristic-domination conclusion
holds at \(p=1\) for \(1<q\le 2\).  The result remains partial relative to the
source problem: \(q>2\) and arbitrary UMD spaces are not settled here.

\section*{Source Problem}

The source is arXiv:1907.11588, PDF page 105, Remark 12.10.  The source states
that the general local-martingale case is open and reduces the missing issue
to an independent-increment comparison.

\begin{center}
\includegraphics[width=\textwidth]{figures/open_problem_crop.png}
\end{center}

The continuation states the aggregate-domination formulation and notes that
even the symmetric case was not known to the author.

\begin{center}

\includegraphics[width=\textwidth]{figures/open_problem_continuation_crop.png}
\end{center}

\section*{Main Result}

For a finite independent family \(x=(x_i)\) of \(L^q(S)\)-valued random
variables, define
\[
  \sfq(x)
  =
  \E\left\|
      \left(\sum_i |x_i|^2\right)^{1/2}
    \right\|_{L^q(S)}.
\]
Let
\[
   \mu_x=\sum_i \mathcal L(x_i)|_{\Lq\setminus\{0\}}
\]
be the aggregate occurrence measure.

\begin{theorem}[Endpoint square-function comparison, \(1<q\le2\)]
\label{thm:square}
Let \(1<q\le2\).  There is an absolute constant \(K\le e/(e-1)\) such that
for every measure space \(S\) and every two finite independent families
\(x=(x_i)\), \(y=(y_j)\) of \(L^q(S)\)-valued random variables,
\[
  \mu_x\le \mu_y
  \quad\Longrightarrow\quad
  \sfq(x)\le K\,\sfq(y).
\]
\end{theorem}

\begin{theorem}[General local \(L^q\), \(p=1\), \(1<q\le2\)]
\label{thm:main}
Let \(1<q\le2\), let \(S\) be a measure space, and put \(X=L^q(S)\).  If
\(N\) and \(M\) are \(X\)-valued local martingales and \(N\) is
characteristically dominated by \(M\) in the sense of
\cite[Definition 12.9]{Yaroslavtsev2019}, then
\[
  \E\sup_{t\ge0}\norm{N_t}_{L^q(S)}
  \le
  C_q
  \E\sup_{t\ge0}\norm{M_t}_{L^q(S)}.
\]
\end{theorem}

\section*{Idea of the Proof}

The point of the endpoint \(p=1\) is that Yaroslavtsev's Banach-function-space
BDG theorem compares martingale maximal functions with lattice square
functions even at \(p=1\) \cite{Yaroslavtsev2018}.  Thus the remaining
question is whether
\[
  \E\left\|\left(\sum_i |x_i|^2\right)^{1/2}\right\|_q
\]
is monotone, up to an absolute constant, under aggregate domination of the
individual jump laws.

For \(1<q\le2\) this square-function norm has exactly the right Laplace
structure.  On a finite weighted \(\ell_q^d\) sublattice, put
\[
  r=q/2,\qquad
  \Phi(v)=\left(\sum_{k=1}^d a_k v_k^r\right)^{1/q},
  \qquad v\in\R_+^d.
\]
Since \(r\le1\), each \(v_k^r\) is the Laplace exponent of a stable
subordinator.  Since \(1/q<1\), the outer power is again a Bernstein
operation.  Therefore
\[
  \Phi(v)=\int_{\R_+^d\setminus\{0\}}
           \left(1-e^{-\langle t,v\rangle}\right)\,\Pi(dt)
\]
for a positive measure \(\Pi\).  For independent positive vectors \(V_i\),
the expectation of the kernel is
\[
  1-\prod_i \E e^{-\langle t,V_i\rangle}.
\]
Writing \(u_i(t)=\E(1-e^{-\langle t,V_i\rangle})\), this quantity is trapped
between constants times \(\min(1,\sum_i u_i(t))\).  The sum
\(\sum_i u_i(t)\) depends only monotonically on the aggregate occurrence
measure.  This is the whole mechanism.

\section*{The Finite-Dimensional Laplace Lemma}

\begin{lemma}[Bernstein representation]
\label{lem:bernstein}
Let \(1<q\le2\), let \(r=q/2\), let \(a_1,\ldots,a_d>0\), and set
\[
  \Phi(v)=\left(\sum_{k=1}^d a_k v_k^r\right)^{1/q},
  \qquad v\in\R_+^d.
\]
There exists a positive measure \(\Pi\) on
\(\R_+^d\setminus\{0\}\) such that, for all \(v\in\R_+^d\),
\[
  \Phi(v)=
  \int_{\R_+^d\setminus\{0\}}
     \left(1-e^{-\langle t,v\rangle}\right)\,\Pi(dt).
\]
\end{lemma}

\begin{proof}
Put \(\alpha=1/q\in(0,1)\) and
\(H(v)=\sum_{k=1}^d a_k v_k^r\).  The scalar identity
\[
  s^\alpha
  =
  c_\alpha\int_0^\infty (1-e^{-\lambda s})\lambda^{-1-\alpha}\,d\lambda,
  \qquad s\ge0,
\]
holds for \(0<\alpha<1\).  Hence
\[
  \Phi(v)
  =
  c_\alpha\int_0^\infty
       \left(1-e^{-\lambda H(v)}\right)
       \lambda^{-1-\alpha}\,d\lambda.
\]
For \(0<r<1\), \(e^{-\lambda a_k v_k^r}\) is the Laplace transform of the
positive \(r\)-stable subordinator at time \(\lambda a_k\); for \(r=1\) it is
the Laplace transform of the point mass at \(\lambda a_k\).  Therefore, for
each \(\lambda>0\), there is a probability measure \(\rho_\lambda\) on
\(\R_+^d\) such that
\[
  e^{-\lambda H(v)}
  =
  \int_{\R_+^d} e^{-\langle t,v\rangle}\,\rho_\lambda(dt).
\]
Thus
\[
  1-e^{-\lambda H(v)}
  =
  \int_{\R_+^d}
      \left(1-e^{-\langle t,v\rangle}\right)\,\rho_\lambda(dt).
\]
Integrating the measures \(c_\alpha\lambda^{-1-\alpha}\rho_\lambda\,d\lambda\)
over \(\lambda>0\) gives the required positive measure \(\Pi\).  The origin
may be discarded because the kernel vanishes at \(t=0\).
\end{proof}

\begin{lemma}[Finite-dimensional aggregate comparison]
\label{lem:finite}
Let \(1<q\le2\).  Let \(V_i,W_j\) be independent \(\R_+^d\)-valued random
vectors.  Suppose that their aggregate occurrence measures satisfy
\[
  \sum_i \Pp(V_i\in A)
  \le
  \sum_j \Pp(W_j\in A),
  \qquad A\in\cB(\R_+^d),\ 0\notin A.
\]
Then, for \(\Phi\) as in Lemma \ref{lem:bernstein},
\[
  \E\Phi\left(\sum_i V_i\right)
  \le
  \frac{e}{e-1}
  \E\Phi\left(\sum_j W_j\right).
\]
\end{lemma}

\begin{proof}
Using Lemma \ref{lem:bernstein} and Tonelli's theorem,
\[
  \E\Phi\left(\sum_i V_i\right)
  =
  \int
  \left[
    1-\prod_i \E e^{-\langle t,V_i\rangle}
  \right]\Pi(dt).
\]
For fixed \(t\), define
\[
  u_i(t)=\E\left(1-e^{-\langle t,V_i\rangle}\right),
  \qquad
  U_V(t)=\sum_i u_i(t),
\]
and define \(U_W(t)\) similarly.  Since \(0\le u_i(t)\le1\),
\[
  1-\prod_i(1-u_i(t))
  \le
  \min(1,U_V(t)).
\]
Also \(\prod_i(1-u_i(t))\le e^{-U_V(t)}\), so
\[
  1-\prod_i(1-u_i(t))
  \ge
  1-e^{-U_V(t)}
  \ge
  (1-e^{-1})\min(1,U_V(t)).
\]
Aggregate domination gives
\[
  U_V(t)
  =
  \int_{\R_+^d\setminus\{0\}}
      (1-e^{-\langle t,v\rangle})\,d\nu_V(v)
  \le
  \int_{\R_+^d\setminus\{0\}}
      (1-e^{-\langle t,v\rangle})\,d\nu_W(v)
  =
  U_W(t).
\]
Consequently the integrand for the \(V_i\)'s is at most
\((1-e^{-1})^{-1}\) times the integrand for the \(W_j\)'s.  Integrating with
respect to \(\Pi\) proves the claim.
\end{proof}

\section*{Proof of the Square-Function Theorem}

\begin{proof}[Proof of Theorem \ref{thm:square}]
First assume that \(L^q(S)\) is a finite weighted \(\ell_q^d\) space with
weights \(a_k>0\).  For \(z\in\ell_q^d\), set
\[
  Tz=(|z_1|^2,\ldots,|z_d|^2)\in\R_+^d.
\]
The aggregate domination \(\mu_x\le\mu_y\) pushes forward under \(T\), so the
positive vectors \(V_i=T x_i\) and \(W_j=T y_j\) satisfy the hypothesis of
Lemma \ref{lem:finite}.  Moreover
\[
  \left\|
      \left(\sum_i |x_i|^2\right)^{1/2}
  \right\|_{\ell_q(a)}
  =
  \Phi\left(\sum_i V_i\right),
\]
with \(\Phi(v)=(\sum_k a_k v_k^{q/2})^{1/q}\).  Lemma \ref{lem:finite}
therefore proves the finite-dimensional claim.

For general \(L^q(S)\), use the standard finite-conditional-expectation
approximation.  Since the finite random families are Bochner measurable, their
essential ranges are separable.  After restricting to the sigma-finite
support generated by a countable dense subset of these ranges, choose positive
finite-rank conditional expectations \(P_n\) onto finite measurable
partitions such that \(P_n z\to z\) in \(L^q(S)\) for every vector in those
essential ranges.

The deterministic map \(P_n:L^q(S)\to P_nL^q(S)\) preserves aggregate
domination:
\[
  \sum_i\Pp(P_nx_i\in B)
  =
  \mu_x(P_n^{-1}B)
  \le
  \mu_y(P_n^{-1}B)
  =
  \sum_j\Pp(P_ny_j\in B)
\]
for every Borel \(B\subset P_nL^q(S)\setminus\{0\}\).  Hence the
finite-dimensional result gives
\[
  \E\left\|
      \left(\sum_i |P_nx_i|^2\right)^{1/2}
    \right\|_q
  \le
  \frac{e}{e-1}
  \E\left\|
      \left(\sum_j |P_ny_j|^2\right)^{1/2}
    \right\|_q.
\]
For every deterministic finite family \((z_i)\),
\[
  \left(\sum_i |P_nz_i|^2\right)^{1/2}
  \le
  P_n\left(\left(\sum_i |z_i|^2\right)^{1/2}\right)
\]
pointwise, by Jensen's inequality for the Hilbert norm.  Thus the right-hand
side is bounded by
\[
  \frac{e}{e-1}
  \E\left\|
      \left(\sum_j |y_j|^2\right)^{1/2}
    \right\|_q.
\]
On the other hand, \(P_nx_i\to x_i\) in \(L^q\) almost surely along the chosen
essential ranges, so the corresponding square functions converge almost
surely in \(L^q\).  Fatou's lemma gives the desired inequality after letting
\(n\to\infty\).
\end{proof}

\section*{Discrete Martingale Consequence}

\begin{proposition}[Discrete \(p=1\) comparison for \(1<q\le2\)]
\label{prop:discrete}
Let \(1<q\le2\).  Let \((\xi_i)_{i=1}^n\) and \((\eta_j)_{j=1}^m\) be
independent mean-zero \(L^q(S)\)-valued random variables.  If
\[
   \sum_i \Pp(\xi_i\in B)
   \le
   \sum_j \Pp(\eta_j\in B)
   \qquad(B\in\cB(\Lq),\;0\notin B),
\]
then
\[
   \E\max_{1\le k\le n}
      \left\|\sum_{i=1}^k \xi_i\right\|_{L^q(S)}
   \le
   C_q
   \E\max_{1\le k\le m}
      \left\|\sum_{j=1}^k \eta_j\right\|_{L^q(S)}.
\]
\end{proposition}

\begin{proof}
The aggregate hypothesis is exactly \(\mu_\xi\le\mu_\eta\).  Therefore
Theorem \ref{thm:square} gives
\[
  \E\left\|\left(\sum_i|\xi_i|^2\right)^{1/2}\right\|_{L^q}
  \le
  \frac{e}{e-1}
  \E\left\|\left(\sum_j|\eta_j|^2\right)^{1/2}\right\|_{L^q}.
\]
Yaroslavtsev's Banach-function-space BDG theorem
\cite[Theorem 4.32 and Proposition 4.33]{Yaroslavtsev2018} says that for
every \(L^q(S)\)-valued martingale difference sequence \(d=(d_i)\),
\[
  \E\sup_N\left\|\sum_{i=1}^N d_i\right\|_{L^q}
  \simeq_q
  \E\left\|
      \left(\sum_i |d_i|^2\right)^{1/2}
    \right\|_{L^q}.
\]
Both independent mean-zero families are martingale difference sequences in
any fixed order.  Applying this equivalence to \((\xi_i)\) and \((\eta_j)\)
proves the maximal comparison.
\end{proof}

\section*{General Local Martingales}

\begin{proof}[Proof of Theorem \ref{thm:main}]
The conditional endpoint packet in this run proves the following reduction:
for a fixed \(1<q<\infty\), if the endpoint square-function aggregate
comparison of Theorem \ref{thm:square} holds in \(L^q(S)\), then the full
\(L^q(S)\)-valued general-local characteristic-domination theorem holds at
\(p=1\).  We recall the argument.

Yaroslavtsev's decoupled tangent construction and canonical decomposition
reduce the comparison to continuous parts and purely discontinuous tangent
jump martingales.  Continuous parts are controlled at \(p=1\) by the source's
continuous characteristic-domination estimate \cite[Proposition 12.6]{Yaroslavtsev2019}.
For the jump part, conditional on the local characteristics and after the
usual truncation of time and jump size, the decoupled tangent jumps are finite
independent mean-zero families.  Characteristic domination controls the total
spatial jump compensator,
\[
  \nu^N(\R_+\times B)\le \nu^M(\R_+\times B),
  \qquad B\in\cB(L^q(S)),\ 0\notin B,
\]
which is exactly the aggregate occurrence-measure domination used in
Proposition \ref{prop:discrete}.

Applying Proposition \ref{prop:discrete} conditionally, then passing through
the truncation and localization limits, gives the \(L^1\) maximal comparison
for the discontinuous tangent martingales.  Combining this with the continuous
part and the \(p=1\) Banach-function-space BDG/canonical-decomposition
equivalence \cite[Section 4.3]{Yaroslavtsev2018} yields the zero-initial-value
case.  If the martingales have nonzero initial values, characteristic
domination gives \(\norm{N_0}_{L^q}\le\norm{M_0}_{L^q}\); applying the
zero-start estimate to \(N-N_0\) and \(M-M_0\) finishes the proof.
\end{proof}

\section*{Why the Argument Stops at \(q=2\)}

The proof is not merely missing a cosmetic extension.  Its central positive
Laplace representation relies on \(r=q/2\le1\).  For \(q>2\), the candidate
\[
  v\mapsto \left(v_1^{q/2}+v_2^{q/2}\right)^{1/q}
\]
is no longer a Bernstein function on the positive cone.  For example at
\(q=4\), \(\partial_{v_1}(v_1^2+v_2^2)^{1/4}\) increases from \(0\) as
\(v_1\) leaves \(0\) with \(v_2>0\) fixed, whereas a Bernstein-function
partial derivative would be completely monotone.  Thus the current proof
cannot simply be pushed past \(q=2\).

Other routes tried during the push also appear to need a new idea for
\(q>2\).  Dualizing
\[
  \left\|\sum_i |x_i|^2\right\|_{q/2}
\]
introduces a random supremum over \(L^{q/(q-2)}\) weights, while scalar
Laplace comparison only handles each fixed weight.  Stable-type
representations of \(\ell_r\) norms for \(1<r\le2\) introduce signed scalar
sums; aggregate domination is not enough to compare signed fractional moments
because cancellation can make the dominating family smaller.  These failures
are recorded as limitations, not as counterexamples.

\section*{Verification and Novelty Notes}

The finite-dimensional proof is a direct positive-kernel calculation and
works in extended expectation values.  The passage to general \(L^q(S)\) uses
only positive finite-rank conditional expectations and Fatou's lemma.

Cheap run indexes contained earlier packets for the \(p>1\) \(L^q\) case,
Hilbert all-\(p\) case, \(\ell_q\) natural-exponent case, \(c_0\)
counterexamples, and the conditional \(p=1\) reduction.
No packet proved the \(L^q(S)\), \(p=1\), \(1<q\le2\) endpoint.  Web searches
for exact phrases around Remark 12.10, \(p=1\), \(L^q\), characteristic
domination, aggregate domination, and square functions did not reveal a later
paper settling this endpoint.

\section*{Human Review Recommendation}

Review Lemmas \ref{lem:bernstein} and \ref{lem:finite} first.  If the
multivariate Bernstein representation is accepted, the square-function
comparison is immediate.  Then check the finite-conditional-expectation
approximation in Theorem \ref{thm:square}; the martingale consequence is the
previous conditional reduction with its conditional square-function lemma now
proved for \(1<q\le2\).

\begin{thebibliography}{9}

\bibitem{Yaroslavtsev2019}
I. S. Yaroslavtsev.
\newblock Local characteristics and tangency of vector-valued martingales.
\newblock arXiv:1907.11588, 2019.

\bibitem{Yaroslavtsev2018}
I. S. Yaroslavtsev.
\newblock Burkholder--Davis--Gundy inequalities in UMD Banach spaces.
\newblock arXiv:1807.05573, 2018.

\bibitem{DirksenYaroslavtsev2017}
S. Dirksen and I. S. Yaroslavtsev.
\newblock \(L^q\)-valued Burkholder--Rosenthal inequalities.
\newblock arXiv:1707.00109, 2017.

\bibitem{SchillingSongVondracek}
R. L. Schilling, R. Song, and Z. Vondracek.
\newblock \emph{Bernstein Functions: Theory and Applications}.
\newblock De Gruyter, 2010.

\bibitem{Feller}
W. Feller.
\newblock \emph{An Introduction to Probability Theory and Its Applications,
  Vol. II}.
\newblock Wiley, 1971.

\end{thebibliography}

\end{document}
